\begin{proposition}[4.7.11]
\label{IV.2.4.7.11-fr}
Sean $k$ un cuerpo, $X$ un preesquema localmente de tipo finito sobre $k$,
$\mathscr F$ un $\mathcal O_X$-módulo coherente, $k'$ una extensión de $k$,
$X'=X\otimes_k k'$, $\mathscr F'=\mathscr F\otimes_k k'$, $x$ un punto de
$X$ y $x'$ un punto de $X'$ por encima de $x$. Para que $\mathscr F$ sea
geométricamente reducido \textup{(resp.} geométricamente puntualmente íntegro\textup{)} en el
punto $x$, es necesario y suficiente que $\mathscr F'$ sea geométricamente
reducido \textup{(resp.} geométricamente puntualmente íntegro\textup{)} en el punto $x'$.
\end{proposition}

\begin{proof}
Las componentes irreducibles de $\operatorname{Supp}(\mathscr F')$ que
contienen a $x'$ dominan cada una una componente irreducible de
$\operatorname{Supp}(\mathscr F)$ que contiene a $x$ y, recíprocamente, cada
una de estas últimas está dominada por una componente irreducible de
$\operatorname{Supp}(\mathscr F')$ que contiene a $x'$ ((4.2.7, (i)) y
(2.3.4)). La afirmación relativa a la propiedad de separabilidad se deduce,
pues, de (4.2.7, (ii)), (4.7.9) y (4.7.10). Sea $Y$ el subpreesquema cerrado
de $X$ definido por el ideal $\mathscr J$ anulador de $\mathscr F$, de modo
que $Y'=Y\otimes_k k'$ es el subpreesquema cerrado de $X'$ definido por el
ideal $\mathscr J'=\mathscr J\otimes_k k'$ anulador de $\mathscr F'$~; decir
que $\mathscr F$ (resp. $\mathscr F'$) es geométricamente puntualmente íntegro
en el punto $x$ (resp. $x'$) equivale a decir que $\mathscr F$ (resp.
$\mathscr F'$) es geométricamente reducido en el punto $x$ (resp. $x'$) y que
$Y$ (resp. $Y'$) es geométricamente puntualmente íntegro en el punto $x$
(resp. $x'$). En virtud de la afirmación del enunciado relativa a la
separabilidad, la hipótesis de que $\mathscr F$ es geométricamente puntualmente
íntegro en el punto $x$, o la hipótesis de que $\mathscr F'$ es
geométricamente puntualmente íntegro en el punto $x'$, implican ambas que
existe en $Y$ una vecindad de $x$ que es separable, y la conclusión se deduce de
(4.6.10).
\end{proof}

\begin{definition}[4.7.12]
\label{IV.2.4.7.12-fr}
Sean $k$ un cuerpo, $X$ un $k$-preesquema localmente de tipo finito sobre $k$,
$\mathscr F$ un $\mathcal O_X$-módulo coherente y $x$ un punto maximal del
soporte de $\mathscr F$. Se llama multiplicidad

\oldpage[IV-2]{80}

total de $x$ (o del
ciclo primo $\overline{\{x\}}$) para $\mathscr F$ (relativa a $k$) al
producto de la multiplicidad radicial de $x$ para $\mathscr F$ por la
multiplicidad separable de $\boldsymbol{k}(x)$ sobre $k$.
\end{definition}

Equivalentemente, la multiplicidad total de $x$ para $\mathscr F$ es el
producto de la \emph{longitud}
$\operatorname{long}_{\mathcal O_x}(\mathscr F_x)$ por la multiplicidad total
de $\boldsymbol{k}(x)$ sobre $k$ (4.7.4). Con las notaciones de (4.7.11), la multiplicidad
total de $x$ para $\mathscr F$ es igual a la \emph{suma} de las multiplicidades
totales para $\mathscr F'$ de los puntos maximales $x'_i$ de
$\operatorname{Supp}(\mathscr F')$ que están por encima de $x$~; además,
existe una extensión \emph{finita} $k'$ de $k$ tal que la multiplicidad total de $x$
para $\mathscr F$ es \emph{igual a}
$\sum_i\operatorname{long}(\mathscr F'_{x'_i})$.

\begin{proposition}[4.7.13]
\label{IV.2.4.7.13-fr}
Bajo las hipótesis y con las notaciones de \textup{(4.7.10)}, para que $\mathscr F$ sea
geométricamente puntualmente íntegro en el punto $x$, basta que $x$ pertenezca
a un solo ciclo primo asociado (necesariamente no inmerso) de $\mathscr F$ y
que la multiplicidad total para $\mathscr F$ del punto genérico $z$ de este
ciclo sea igual a $1$.
\end{proposition}

\begin{proof}
En efecto, si $k'$ es una extensión finita de $k$, $x'$ un punto de
$X'=X\otimes_k k'$ por encima de $x$ y
$\mathscr F'=\mathscr F\otimes_k k'$, $x'$ solo puede pertenecer a un ciclo
primo asociado a $\mathscr F'$, pues, por hipótesis, solo hay un punto maximal
de $\operatorname{Supp}(\mathscr F')$ por encima de $z$~; además, $\mathscr F$
es geométricamente reducido por (4.7.10), de donde se sigue la conclusión.
\end{proof}
